\section*{Introduction générale}

Audit Énergie est une entreprise initialement positionnée sur le courtage en énergie.
Son activité historique consiste à accompagner des professionnels dans la renégociation
de leurs contrats, dans un marché devenu progressivement plus concurrentiel et plus
complexe. Cette activité a permis à l'entreprise de développer une connaissance fine de
la relation client, de la prospection commerciale et des compétences nécessaires aux
métiers de la vente à distance.

Depuis plusieurs années, cette expérience métier a progressivement conduit l'entreprise
à se tourner vers la formation professionnelle. Dans un premier temps, cette activité
répondait principalement à un besoin interne : former des profils capables d'intégrer les
équipes commerciales ou de téléprospection. Plus récemment, Audit Énergie a cherché à
structurer cette activité comme un véritable organisme de formation externe, capable
d'accompagner des apprenants dans la préparation de titres professionnels, notamment
dans les domaines de la relation client à distance ou de l'employé commercial.

Ce changement d'échelle a fait apparaître une contrainte centrale. Pour une petite
structure, la formation représente une opportunité de développement, mais aussi un coût
important. Le modèle le plus naturel aurait été de faire intervenir des formateurs
freelance pour dispenser les cours en ligne. Cette solution présente des avantages
évidents : elle repose sur une interaction humaine directe, permet de répondre aux
questions des élèves en temps réel et garantit une certaine souplesse pédagogique.
Cependant, son coût rendait les marges trop faibles pour que le modèle soit réellement
rentable à grande échelle.

Une première solution a donc été de remplacer le cours en direct par un cours audio
préparé à l'avance. Pour répondre à ce besoin, j'ai développé une plateforme capable
d'héberger ces audios et de les diffuser automatiquement le jour de la formation. Le
travail était alors séparé en deux tâches : la rédaction du contenu, éventuellement aidée
par l'intelligence artificielle, puis l'enregistrement de sa lecture. Cette organisation
réduisait le coût par rapport à l'intervention d'un formateur freelance en direct, tout en
permettant aux élèves de suivre la formation selon le calendrier prévu.

Cette approche constituait une première forme d'optimisation, mais elle ne réglait pas
tous les problèmes. La production des cours restait manuelle, chronophage et dépendante
de plusieurs personnes. Des essais de synthèse vocale ont ensuite été réalisés pour
automatiser la lecture des cours. L'idée était simple : si le contenu pouvait être écrit à
l'avance, alors une voix générée par IA pouvait peut-être remplacer l'enregistrement
humain. Mais les premières solutions testées présentaient elles aussi des limites. Les
outils les plus qualitatifs étaient coûteux lorsqu'il fallait produire plusieurs heures
d'audio, tandis que les solutions plus économiques donnaient parfois un résultat moins
adapté à une formation longue.

Cette étape a été importante, car elle a déplacé la réflexion. Le problème n'était pas
seulement de trouver une voix artificielle capable de lire un texte. Il fallait concevoir une
chaîne complète capable de produire un cours exploitable : un texte structuré, oral,
calibré en durée, conforme au programme, suffisamment naturel pour être écouté
plusieurs heures, puis transformé en audio et intégré à la plateforme. Une bonne voix ne
suffit pas si le texte lu est répétitif, mal construit ou trop long pour le temps disponible.

Une autre limite est rapidement apparue : un cours enregistré, même de bonne qualité,
ne remplace pas entièrement un formateur. Lorsqu'un apprenant ne comprend pas un
passage, il doit pouvoir poser une question. Or, dans un modèle fondé uniquement sur
des audios préenregistrés, aucun humain n'est disponible pour répondre en direct. Si
l'objectif était réellement de construire une plateforme de formation autonome, il fallait
donc traiter deux dimensions complémentaires : automatiser la production et la diffusion
des contenus, mais aussi fournir un accompagnement pédagogique capable de répondre
aux questions des élèves de manière contextualisée.

C'est dans ce contexte qu'a émergé le projet étudié dans ce mémoire : concevoir une
plateforme de formation autonome alimentée par l'intelligence artificielle. Cette
plateforme ne devait pas seulement diffuser des fichiers audio. Elle devait être capable de
générer des contenus de formation, de les transformer en audio, de produire des supports
visuels, de proposer des exercices et d'accompagner les élèves grâce à un assistant
conversationnel s'appuyant sur les documents pédagogiques.

La problématique de ce mémoire peut donc être formulée ainsi :

\begin{quote}
Comment concevoir une plateforme de formation autonome alimentée par
l'intelligence artificielle, capable de produire, diffuser et accompagner des contenus
pédagogiques, tout en garantissant leur qualité, leur traçabilité et leur pertinence pour
les apprenants ?
\end{quote}

Cette problématique soulève plusieurs objectifs. Le premier est d'automatiser la
production des cours, en passant d'un référentiel ou d'un programme de formation à des
journées pédagogiques complètes. Le deuxième est de transformer ces contenus en
audios exploitables, avec un coût inférieur à une production humaine répétée et une
qualité suffisante pour un usage prolongé. Le troisième est d'enrichir l'expérience de
formation par des supports visuels et des exercices générés automatiquement. Le
quatrième est de permettre aux apprenants de poser des questions à un assistant
pédagogique capable de répondre à partir du contenu réel de la formation, et non comme
un chatbot généraliste.

Les contributions de ce travail s'organisent donc autour de deux axes principaux. Le
premier concerne la conception d'une pipeline de génération de contenus pédagogiques.
Cette pipeline a progressivement évolué d'une génération directe de texte vers une
chaîne plus structurée, fondée sur un plan pédagogique, des étapes de contrôle qualité,
des artefacts intermédiaires et une meilleure traçabilité des corrections. Le second axe
concerne l'accompagnement des apprenants, notamment par la mise en place d'un
système de questions-réponses contextualisé de type RAG, destiné à réduire les limites
d'un modèle de langage utilisé seul.

La méthode suivie repose sur une démarche itérative. Plusieurs solutions ont été
comparées ou testées : intervention humaine, enregistrement différé, synthèse vocale
externe, génération directe par modèle de langage, puis pipeline plus structurée. Chaque
étape a permis d'identifier de nouvelles contraintes : coût, qualité audio, longueur des
textes, cohérence pédagogique, conformité, capacité à répondre aux élèves et besoin
d'évaluer les résultats. Cette progression constitue le fil conducteur du mémoire : il ne
s'agit pas seulement d'ajouter de l'intelligence artificielle à une plateforme existante,
mais de comprendre comment encadrer cette automatisation pour la rendre fiable,
contrôlable et mesurable.

Enfin, ce travail ne peut pas être évalué uniquement par le fait que la plateforme
fonctionne techniquement. Il doit aussi montrer ce que l'automatisation apporte. Les
résultats seront donc discutés à travers plusieurs types de métriques : réduction des
interventions humaines, coût et temps de production des contenus, qualité des audios,
cohérence pédagogique, satisfaction des apprenants, résultats aux exercices, ainsi que
métriques propres à l'assistant RAG, comme la pertinence des passages retrouvés et la
fidélité des réponses produites.

Ce mémoire est structuré en trois parties. La première présente le contexte de
l'entreprise, les limites du fonctionnement initial et la problématique générale de la
formation autonome. La deuxième étudie les approches existantes : plateformes de
formation en ligne, synthèse vocale, génération de contenus pédagogiques par IA et
systèmes de questions-réponses augmentés par récupération d'information. La troisième
partie présente la solution développée, ses choix d'architecture, les difficultés rencontrées,
les arbitrages réalisés et les méthodes envisagées pour évaluer ses apports.
